\subsection{Sets and Boolean Logic}

\content{
        \begin{definition} A \textbf{set} is a collection of distict 
        objects, called elements of the set.

        A set can be defined by describing it's contents, or listing
        the elements of the set, in curly brackets. 
        \end{definition}
}{% presentation
}{% nopresentation
}{% comments
Describe how we will typically give a set a name, say $A$, and then we
can make statements like $x\in A$. Define $\in$. Define $\emptyset$.
}

\content{
        \begin{definition} The \textbf{union} of two sets $A$ and $B$
        is the set consisting of all elements of $A$ or $B$. The
        union of $A$ and $B$ is denoted $A\cup B$.
        \end{definition}

        \begin{definition} The \textbf{intersection} of two sets $A$
        and $B$ is the set consisting of all elements which belong to
        both $A$ and $B$, denoted $A\cap B$.
        \end{definition}
}{% presentation
}{% nopresentation
}{% comments
}

\content{
        \begin{example} Find the union of the two given sets. $A$ is
        the set of all even numbers. $B$ is the set of whole numbers
        between 1 and 11 (including 1 and 11). 
        \end{example}
}{% presentation
\vspace{1in}
}{% nopresentation
\vspace{1in}
}{% comments
}

\content{
        \begin{definition} A \textbf{subset} of a set $A$ is another
        set $B$, all of whose elements are in $A$. Note that $B$ does
        not have to include \textit{all} the elements of $A$.
        \end{definition}

        \begin{definition} The \textbf{compliment} of a set $A$ is the
        set of elements which do not belong to $A$ (here $A$ must be a
        subset of some set $B$). This is denoted $A^c$ or $B\setminus 
        A$.
        \end{definition}

}{% presentation
}{% nopresentation
}{% comments
}

\content{
        \begin{example} Is $\{1,3,5,7,9\}$ a subset of
        $\{5,7,1,9,2,3,-3\}$?
        \end{example}
}{% presentation
\vspace{1in}
}{% nopresentation
\vspace{1in}
}{% comments
}

\content{
        \begin{example} The set $A=\{1,3,5,7,9\}$ is a subset of
        $B=\{1,2,3,4,5,6,7,8,9,10\}$, find $B\setminus A$.
        \end{example}
}{% presentation
\vspace{1in}
}{% nopresentation
\vspace{1in}
}{% comments
}

\content{
        \begin{definition} A \textbf{statement} is either true or false.
        \end{definition}
        \begin{definition} Boolean Logic combines multiple statements
        that are either true or false into an expression that is
        either true or false. 
        \end{definition}
}{% presentation
\vspace{2in}
}{% nopresentation
\vspace{2in}
}{% comments
Mention that three of the most common operations in boolean logic are 'and`,
'or`, and 'not`, and that these correspond to three set operations.
}

\content{
        \begin{example} Describe the numbers which meet the condition: "odd
        number and less than 20 and greater than 0 and (multiple of 3 or
        multiple of 5)
        \end{example}
}{% presentation
\vspace{1.5in}
}{% nopresentation
\vspace{1.5in}
}{% comments
}

\content{
        \begin{definition} A \textbf{conditional} is a compound
        statement of the form "if $p$, then $q$`` or 
        "if $p$, then $q$, else $r$``.
        \end{definition}
}{% presentation
}{% nopresentation
}{% comments
}

\content{
        \begin{example} In common language, we might say something
        like the following: 
        "If it's raining, we'll go to the mall. Otherwise, we will go
        on a hike.``
        \end{example}
}{% presentation
}{% nopresentation
}{% comments
}

\content{
        \begin{definition} A \textbf{universial quantifier} states that an
        entire set shares a characteristic. 
        \end{definition}

        \begin{definition} A \textbf{existential quantifier} states that a set
        contains at least one element.
        \end{definition}
}{% presentation
}{% nopresentation
}{% comments
Something interesting happens when you negate a quantifier...
}

\content{
        \begin{example} Somebody likes baseball.
        \end{example}
}{% presentation
}{% nopresentation
}{% comments
}

\content{
        \begin{example} Suppose your friend says "Everybody cheats on their
        taxes.`` What is the minimum amount of evidence you need to prove your
        friend wrong? 
        \end{example}

}{% presentation
}{% nopresentation
}{% comments
}
